
% ------------------------------------------------------------
\section{Conclusions}
\label{sec:conclusions}
% ------------------------------------------------------------

PDE and numerics tools have been used from the course, i.e., Fokker--Planck equations, adjoint
(backward) formulations, integration by parts and self-adjointness, conservative
finite-difference discretisation, Crank--Nicolson time stepping, to analyse
mean first passage times in a spatially heterogeneous 1D environment.  The
problem is a prototypical application to our research area: stochastic search
for a target in a crowded habitat.

The main analytical contributions are:
\begin{itemize}[leftmargin=1.3em, noitemsep]
\item a derivation of both It\^o and Fick forms of the FPE
(Section~\ref{sec:derivation});
\item the adjoint/backward derivation for Fick's form on a periodic ring with
a $\delta$-type trap, including the flux-jump condition \eqref{eq:robin} and
the solvability constant $C=\pi$ (Section~\ref{sec:backward});
\item the general additive decomposition
$\langle T\rangle=2\pi/\kappa+\pi^{-1}\!\int(\pi-y)^2/D(y)\,dy$
\eqref{eq:T-avg-general};
\item explicit closed forms for $T(x)$ and $\langle T\rangle$ on the tent
profile, \eqref{eq:T-tent} and \eqref{eq:T-avg-tent}, with all integrals in
elementary form.
\end{itemize}

The main numerical contributions are threefold verification of the analytical
result:
\begin{itemize}[leftmargin=1.3em, noitemsep, label=$\diamond$]
\item finite-difference solution of the backward equation, with empirical
second-order convergence in $h$;
\item Crank--Nicolson finite-difference time evolution of the forward FPE,
with MFPT recovered via survival probabilities;
\item Monte Carlo simulation of the It\^o SDE with a Poisson trap rule.
\end{itemize}
All three methods match the analytical closed form within their expected
errors, demonstrating the connection between PDE, stochastic-process, and
particle-level formulations.

The principal qualitative insight is the role of the relationship between the
trap strength $\kappa$ and the local diffusivity $D(0)$.  With $\kappa$
constant, $\langle T\rangle$ is monotone in the tent parameter $D_1$, so there
is no interior optimum and the additive decomposition \eqref{eq:T-avg-general}
makes the trap contribution and the transport contribution independent.  When
$\kappa$ depends on $D(0)$ the two contributions couple; for the biologically
natural choice $\kappa=\kappa_0/D(0)$ a non-trivial optimum $D_1^*\approx -0.45$
emerges, balancing the advantage of faster transport (large $D_1$) against
the advantage of shorter dwell time at a slow-motion target (small $D_1$).
This is a tidy, quantitative illustration of the speed--accuracy tradeoff that
is thought to underlie many biological search
strategies~\cite{Benichou2011,Peleg2018}.

Natural extensions include: (i) non-symmetric diffusivity profiles (lifted
via scale functions and speed measures~\cite{Pavliotis2014,Gardiner2009});
(ii) noise-induced drift from asymmetric jump rates, giving
$D(x)T''+\mu(x)T'=-1$ (taxis-like behaviour toward the target); (iii)
subdiffusive regimes governed by fractional Fokker--Planck equations,
for which $\langle T\rangle$ diverges~\cite{Metzler2000}; and (iv) the
genuinely two- and three-dimensional problem with radially symmetric $D(r)$,
relevant to nest geometries.
